Personalized AI assistants face a persistent failure: flat memory systems retain stale facts after a user's life changes, causing outdated recommendations. We address this with \textbf{cascading memory invalidation}, a mechanism that propagates invalidation signals through typed dependency edges in a user memory graph without requiring explicit user retraction.

Our approach represents each user's memories as nodes in a directed graph with two edge types: \locatedin edges (linking location-dependent memories to a root location node) and \conflictswith edges (directed from a current preference to an outdated one). When a root node changes, the cascade automatically downweights all dependent memories proportional to edge strength and propagation depth.

We evaluate on two benchmarks. On \locomo (35 real conversational dialogues), structural cascade via \locatedin edges achieves \textbf{81.7\% accuracy} in identifying and invalidating location-dependent memories after a user relocation event, exceeding the 80\% target. We also show that \gptfouromini can establish \conflictswith edges at memory write time with \textbf{100\% precision and 90.9\% recall} on a 19-pair conflict benchmark, solving the \textit{semantic bridging problem} of inferring that ``prefers quiet'' conflicts with ``likes bars.'' On \horizonbench (4,245 preference evolution items), our full transitive cascade improves evolved-preference accuracy from 40.1\% (\flatmemory baseline) to \textbf{72.2\%}, a \textbf{+32.2 percentage point} gain. Crucially, we identify \textit{bidirectional edge cancellation} as a failure mode: using undirected \conflictswith edges collapses performance back to the flat memory baseline. These results establish directed graph cascades as a practical and effective approach to automated memory maintenance in personalized AI systems.
